\chapter{Marge, liquidation et répartition des pertes}\label{ch:m3:margin-liquidation-and-loss-allocation}

Le 10 octobre 2025, après l'annonce d'un droit de douane de 100~\% sur les importations chinoises aux États-Unis,
environ 19~milliards de dollars de positions crypto à effet de levier ont été liquidées en vingt-quatre heures, et les positions ouvertes
sur \omterm{def:m3:perpetual-futures:perp}{contrats perpétuels} dans les principales plateformes ont chuté de 43~\%, passant de 217~milliards de dollars à 123~milliards. Le bitcoin a reculé
jusqu'à environ 106\,500~dollars, et les \omterm{def:m3:blockchains-for-traders:key}{jetons} suivis par CoinDesk Research ont perdu environ 47~\% en moyenne. Sur certaines plateformes, les
liquidations ont été si importantes et si rapides que les coussins destinés à absorber leurs pertes ont été épuisés, et les
plateformes ont clôturé les positions bénéficiaires d'autres traders, qui n'avaient rien fait de mal, pour équilibrer leurs comptes.
Une plateforme de produits dérivés sur crypto-actifs n'a pas de membres compensateurs entre elle et ses clients, et aucun
fonds de défaillance alimenté par des banques~: la marge propre de ses clients, son \omterm{def:m3:margin-liquidation-and-loss-allocation:engine}{moteur de liquidation}, son \omterm{def:m3:margin-liquidation-and-loss-allocation:fund}{fonds d'assurance} et, en dernier
ressort, ses clients gagnants constituent la totalité de la cascade de défaillance. Ce chapitre suit une position à effet de levier
depuis sa marge jusqu'à son \omterm{def:m3:margin-liquidation-and-loss-allocation:prices}{prix de liquidation}, à travers le moteur qui la clôture, puis le long de la cascade~; il construit ensuite
une cascade et en fait l'analyse.

\section{Marge isolée, marge croisée et marge en portefeuille}

Sur une bourse de contrats à terme réglementée, la marge d'un client est fixée par la chambre de compensation et collectée par
l'intermédiaire d'un membre compensateur (livre~1, chapitre~20). Une plateforme crypto fixe et collecte elle-même la marge, sur
chaque compte, dans le collatéral propre au compte, en continu.

\begin{definition}[Marge isolée, marge croisée]\label{def:m3:margin-liquidation-and-loss-allocation:modes}
Avec une \emph{marge isolée}\index{marge isolée}, chaque position dispose de sa propre marge, assignée par le trader~;
ses pertes ne peuvent consommer que cette marge, et elle n'est liquidée que seule. Avec une \emph{marge
croisée}\index{marge croisée}, toutes les positions d'un compte partagent le collatéral du compte~; les gains sur une position compensent
les pertes sur une autre, et le compte est liquidé lorsque ses fonds propres totaux passent sous la somme de ses marges
de maintenance.
\end{definition}

La \omterm{def:m3:margin-liquidation-and-loss-allocation:modes}{marge isolée} limite ce qu'un seul pari peut perdre et permet de calculer facilement son \omterm{def:m3:margin-liquidation-and-loss-allocation:prices}{prix de liquidation}~; la \omterm{def:m3:margin-liquidation-and-loss-allocation:modes}{marge croisée}
utilise le collatéral de façon plus efficace et résiste à une perte sur une position si les autres gagnent, au prix du risque de perdre
tout le contenu du compte d'un seul coup. La marge en portefeuille (livre~1, chapitre~20), proposée par certaines plateformes
aux gros comptes, fixe l'exigence à partir de la perte en scénario de l'ensemble du portefeuille, y compris les avoirs au comptant
utilisés comme collatéral. La marge de maintenance est généralement échelonnée~: le taux augmente avec la taille de la position, car
une position plus grande prend plus de temps et coûte plus cher à liquider.

\begin{dated}{2026-09}{Marge de maintenance échelonnée}\label{dat:m3:margin-liquidation-and-loss-allocation:tiers}
Binance calcule la marge de maintenance comme la valeur notionnelle de la position multipliée par le taux de marge de maintenance de sa
tranche, diminuée d'un montant de maintenance qui rend l'exigence continue d'une tranche à la suivante~;
les tranches, les taux et l'effet de levier maximal sont publiés pour chaque contrat et ajustés. Les données d'instrument de BitMEX
au 24 septembre 2026 donnaient pour XBTUSD une marge initiale de 1~\% et une marge de maintenance de 0.5~\% au premier palier de limite de risque (base risk limit).
\end{dated}

\section{Le prix de liquidation}

\begin{definition}[Prix de liquidation, prix de faillite]\label{def:m3:margin-liquidation-and-loss-allocation:prices}
Le \emph{prix de liquidation}\index{prix de liquidation} d'une position est le \omterm{def:m3:perpetual-futures:index}{prix de marquage} auquel sa marge plus le
P\&L latent tombe au niveau de sa marge de maintenance, si bien que la plateforme la reprend. Son \emph{prix de
faillite}\index{prix de faillite} est le prix auquel la marge plus le P\&L latent sont exactement nuls.
\end{definition}

\begin{proposition}[Prix de liquidation des positions isolées]\label{prop:m3:margin-liquidation-and-loss-allocation:liq}
Soit une position linéaire de $q$ unités ouverte au prix $E$, avec une marge $M$ et un taux de marge de maintenance $m$ (sans
montant de maintenance). Une position longue ($q > 0$) est liquidée et fait faillite aux prix
\[
P_{\mathrm{liq}} = \frac{qE - M}{q(1 - m)}, \qquad P_{\mathrm{bk}} = E - \frac{M}{q},
\]
et une position courte de $n = -q$ unités est liquidée et fait faillite aux prix $P_{\mathrm{liq}} = (M + nE)/(n(1 + m))$ et $P_{\mathrm{bk}} = E + M/n$. Avec
un effet de levier $\Lambda = qE/M$, le \omterm{def:m3:margin-liquidation-and-loss-allocation:prices}{prix de liquidation} de la position longue vaut $E(1 - 1/\Lambda)/(1 - m)$. Une position longue inverse de
$N$ contrats d'un dollar chacun, avec une marge $M$ en coin, est liquidée au prix $N(1 + m)/(M + N/E)$.
\end{proposition}

\begin{proof}
Pour la position longue, $M + q(P - E) = m\,qP$ donne la première formule et $M + q(P - E) = 0$ la seconde~; celles de la position courte
sont identiques, au signe du P\&L près. Pour la position longue inverse, les fonds propres en coin valent $M + N(1/E - 1/P)$ et
l'exigence de maintenance $mN/P$~; en les égalant, on obtient le prix énoncé.
\end{proof}

Le symbole $\Lambda$ n'est utilisé que dans ce chapitre (dans le livre~4, il désigne un compensateur). Avec un effet de levier de 10, une
position longue est liquidée environ 9.5~\% sous son prix d'entrée, et à un effet de levier de 50, environ 1.5~\% en dessous. La position longue inverse est liquidée plus tôt
que la position linéaire à même effet de levier, car sa marge en coin perd de la valeur dans le même mouvement~: elle est longue deux fois, comme
le \cref{ch:m3:perpetual-futures} l'a expliqué. \Cref{fig:m3:margin-liquidation-and-loss-allocation:liqprice} représente les trois cas.

\begin{omfigure}
\begin{tikzpicture}
\begin{axis}[width=11cm,height=5cm,xmode=log,log basis x=10,
  xlabel={effet de levier},ylabel={prix de liquidation, \% du prix d'entrée},
  tick label style={font=\small},label style={font=\small},
  xmin=2,xmax=100,ymin=40,ymax=160,grid=major,grid style={gray!25},
  xtick={2,5,10,20,50,100},xticklabels={2,5,10,20,50,100},
  legend columns=3,legend style={font=\footnotesize,at={(0.5,-0.3)},anchor=north,draw=none,fill=none,/tikz/every even column/.append style={column sep=8pt}},
  table/col sep=comma]
\addplot[omThm,very thick] table[x=leverage,y=long_linear]{figdata/markets-3/18-margin-liquidation-and-loss-allocation/liqprice.csv};\addlegendentry{longue, linéaire}
\addplot[omDef,very thick,dashed] table[x=leverage,y=short_linear]{figdata/markets-3/18-margin-liquidation-and-loss-allocation/liqprice.csv};\addlegendentry{courte, linéaire}
\addplot[omExo,very thick,dotted] table[x=leverage,y=long_inverse]{figdata/markets-3/18-margin-liquidation-and-loss-allocation/liqprice.csv};\addlegendentry{longue, inverse}
\end{axis}
\end{tikzpicture}

\omcaption{\omterm{def:m3:margin-liquidation-and-loss-allocation:prices}{Prix de liquidation} des positions isolées en fonction de l'effet de levier, avec un taux de marge de maintenance de
0.5~\% (\cref{prop:m3:margin-liquidation-and-loss-allocation:liq}). Une position longue inverse avec marge en coin est
liquidée plus tôt qu'une position linéaire~: à un effet de levier de 2, après une baisse de 33~\% au lieu de 50~\%. Données~: le tutoriel du
chapitre.}\label{fig:m3:margin-liquidation-and-loss-allocation:liqprice}
\end{omfigure}

Deux conséquences comptent pour le risque. La liquidation est déclenchée par le \omterm{def:m3:perpetual-futures:index}{prix de marquage} (\cref{ch:m3:perpetual-futures}), si bien que le
\omterm{def:m3:margin-liquidation-and-loss-allocation:prices}{prix de liquidation} d'un trader dépend de l'indice de la plateforme, et non des transactions qu'il voit~; et l'écart entre le \omterm{def:m3:margin-liquidation-and-loss-allocation:prices}{prix de liquidation}
et le \omterm{def:m3:margin-liquidation-and-loss-allocation:prices}{prix de faillite}, c'est-à-dire la marge de maintenance, est tout ce dont dispose la plateforme pour clore la position sans perte.

\section{Moteurs de liquidation}

\begin{definition}[Moteur de liquidation]\label{def:m3:margin-liquidation-and-loss-allocation:engine}
Un \emph{moteur de liquidation}\index{moteur de liquidation} est le système de la plateforme qui surveille chaque compte au regard de
sa marge de maintenance au \omterm{def:m3:perpetual-futures:index}{prix de marquage}, reprend les positions des comptes qui la franchissent et les clôture,
en général en envoyant des ordres dans le carnet d'ordres de la plateforme, parfois en les transférant à des
fournisseurs de liquidité désignés.
\end{definition}

Le moteur est la chambre de compensation de la plateforme, sous forme logicielle. Il annule les ordres ouverts du compte, peut d'abord tenter
une liquidation partielle qui ramène la position à une tranche de marge inférieure, puis vend (pour une position longue) dans les cours acheteurs. S'il vend
au-dessus du \omterm{def:m3:margin-liquidation-and-loss-allocation:prices}{prix de faillite}, ce qui reste de la marge du trader est conservé, sur beaucoup de plateformes, par le \omterm{def:m3:margin-liquidation-and-loss-allocation:fund}{fonds d'assurance}~;
s'il vend en dessous, le fonds verse la différence. Le moteur est rapide, mécanique et indifférent au prix~: ses ordres
sont exactement les ventes forcées des spirales de liquidité décrites au livre~1, chapitre~31, qui arrivent lorsque le carnet est le plus mince.
Les pages d'aide de FTX décrivaient l'alternative~: des fournisseurs de liquidité de secours, invités parmi ses paliers de volume les plus élevés, gardaient au moins 500\,000~dollars
sur la plateforme et s'engageaient à absorber au moins 0.1~million de liquidations par minute et 0.3~million par heure, en reprenant la position
et le collatéral d'un compte avant sa faillite. Et la plateforme elle-même peut faillir au moment où on en a besoin~: le 13 mars 2020, BitMEX a subi
deux attaques par déni de service, à 02:16 et 12:56 UTC, qui ont retardé ou empêché les requêtes adressées à la plateforme.

\section{Fonds d'assurance, auto-deleveraging et pertes mutualisées}

\begin{definition}[Fonds d'assurance]\label{def:m3:margin-liquidation-and-loss-allocation:fund}
Un \emph{fonds d'assurance}\index{fonds d'assurance} est un ensemble d'actifs que détient une plateforme de produits dérivés pour absorber les pertes
des positions liquidées sous leurs \omterm{def:m3:margin-liquidation-and-loss-allocation:prices}{prix de faillite}, alimenté surtout par la marge restante des positions
liquidées au-dessus de ces prix.
\end{definition}

\begin{dated}{2026-09}{Solde d'un fonds d'assurance}\label{dat:m3:margin-liquidation-and-loss-allocation:fund}
Le point d'accès public aux données d'assurance de BitMEX indiquait, au 23 septembre 2026, un \omterm{def:m3:margin-liquidation-and-loss-allocation:fund}{fonds d'assurance} de quelque 3\,700
bitcoins et un fonds en USDT d'environ 30.6 millions d'USDT. Les plateformes publient le solde de leurs fonds de différentes façons et à des
fréquences différentes.
\end{dated}

\begin{definition}[Auto-deleveraging, perte mutualisée]\label{def:m3:margin-liquidation-and-loss-allocation:adl}
L'\emph{auto-deleveraging}\index{auto-deleveraging} (ADL) est la clôture de positions opposées rentables d'autres traders,
au \omterm{def:m3:margin-liquidation-and-loss-allocation:prices}{prix de faillite} de la position en défaut ou au \omterm{def:m3:perpetual-futures:index}{prix de marquage}, lorsque le \omterm{def:m3:margin-liquidation-and-loss-allocation:fund}{fonds d'assurance} ne peut pas absorber la perte liée à une
liquidation, les traders étant sélectionnés selon un classement publié. Une \emph{perte mutualisée}\index{perte mutualisée} est une
perte répartie sur un groupe de traders, par exemple sur l'ensemble des comptes rentables au prorata de leurs gains, plutôt que
d'être imposée à ceux sélectionnés par un classement.
\end{definition}

\begin{omfigure}
\centering
\begin{tikzpicture}[font=\small,
  w/.style={draw,thick,rounded corners=2pt,align=center,minimum height=0.85cm,text width=11.5cm},
  arr/.style={-{Stealth[length=2.2mm]},thick}]
\node[w,draw=omDef,fill=omDef!8] (a) at (0,4.4) {1.\ la marge propre du trader, jusqu'au prix de faillite};
\node[w,draw=omThm,fill=omThm!8] (b) at (0,3.0) {2.\ le moteur de liquidation vend dans le carnet~: excédent ou déficit};
\node[w,draw=omProp,fill=omProp!8] (c) at (0,1.6) {3.\ le fonds d'assurance couvre le déficit};
\node[w,draw=omExo,fill=omExo!8] (d) at (0,0.2) {4.\ auto-deleveraging~: positions opposées rentables clôturées, selon leur rang};
\node[w,draw=gray,fill=gray!8] (e) at (0,-1.2) {5.\ (sur certaines plateformes, ou historiquement) pertes mutualisées entre les gagnants};
\draw[arr] (a) -- (b); \draw[arr] (b) -- (c); \draw[arr] (c) -- (d); \draw[arr] (d) -- (e);
\end{tikzpicture}

\omcaption{La cascade de défaillance d'une plateforme de produits dérivés sur crypto-actifs. Il n'y a ni membres compensateurs ni fonds de défaillance
alimenté par eux, comme dans le livre~1, chapitre~5~: après la marge du perdant et le \omterm{def:m3:margin-liquidation-and-loss-allocation:fund}{fonds d'assurance} de la plateforme, les
pertes retombent sur les gagnants. Schématique.}\label{fig:m3:margin-liquidation-and-loss-allocation:waterfall}
\end{omfigure}

\begin{dated}{2026-09}{Règles d'auto-deleveraging sur deux plateformes}\label{dat:m3:margin-liquidation-and-loss-allocation:adlrules}
Binance n'applique l'\omterm{def:m3:margin-liquidation-and-loss-allocation:adl}{auto-deleveraging} que si ses \omterm{def:m3:margin-liquidation-and-loss-allocation:fund}{fonds d'assurance} ne peuvent pas absorber une position en faillite~; les positions
opposées rentables sont classées selon le pourcentage de P\&L multiplié par le levier effectif et clôturées au \omterm{def:m3:margin-liquidation-and-loss-allocation:prices}{prix de faillite} de
l'ordre liquidé, sans frais de négociation, et chaque position affiche un indicateur de son rang dans la file d'attente. La documentation de Hyperliquid
déclenche l'\omterm{def:m3:margin-liquidation-and-loss-allocation:adl}{auto-deleveraging} lorsque la valeur d'un compte ou d'une position isolée devient négative, classe
les utilisateurs rentables selon (\omterm{def:m3:perpetual-futures:index}{prix de marquage} sur prix d'entrée) multiplié par (notionnel sur valeur du compte), et clôture leurs positions au
\omterm{def:m3:perpetual-futures:index}{prix de marquage} précédent contre l'utilisateur en déficit.
\end{dated}

Le classement désigne qui paie~: les gagnants les plus rentables et les plus fortement endettés passent en premier, car leur clôture retire
le plus d'exposition pour le plus petit nombre de comptes. Pour le gagnant, l'\omterm{def:m3:margin-liquidation-and-loss-allocation:adl}{auto-deleveraging} est une sortie involontaire à un prix
qu'il n'a pas choisi, souvent au milieu du mouvement sur lequel il était positionné, et ses couvertures ailleurs restent ouvertes. Un
teneur de marché vendeur du \omterm{def:m3:perpetual-futures:perp}{contrat perpétuel} contre le spot peut voir sa vente clôturée et sa position longue au spot laissée à nu lors d'un krach. Après
octobre 2025, une analyse de Tarun Chitra (Gauntlet), relayée par Protos, soutenait que le classement par file d'attente, commun à
de nombreuses plateformes avait clôturé environ 650~millions de dollars de positions rentables de plus que le minimum nécessaire sur Hyperliquid, et
proposait à la place une règle de répartition au prorata tenant compte du risque.

\section{Une cascade de liquidations disséquée}

\begin{definition}[Cascade de liquidations]\label{def:m3:margin-liquidation-and-loss-allocation:cascade}
Une \emph{cascade de liquidations}\index{cascade de liquidations} est une suite dans laquelle les ventes forcées des \omterm{def:m3:margin-liquidation-and-loss-allocation:engine}{moteurs de liquidation}
font bouger le prix assez pour atteindre les \omterm{def:m3:margin-liquidation-and-loss-allocation:prices}{prix de liquidation} d'autres positions, dont les ventes forcées le font encore bouger.
\end{definition}

Qu'un choc de prix s'arrête ou se propage en cascade dépend du nombre d'unités situées à chaque \omterm{def:m3:margin-liquidation-and-loss-allocation:prices}{prix de liquidation} et de la distance
dont chaque unité vendue fait bouger le prix.

\begin{proposition}[Le multiplicateur local de cascade]\label{prop:m3:margin-liquidation-and-loss-allocation:kappa}
Supposons que chaque unité vendue fasse baisser le prix de $k$ (impact linéaire et permanent), et notons $\rho(P)$ le nombre d'unités
dont les \omterm{def:m3:margin-liquidation-and-loss-allocation:prices}{prix de liquidation} se situent dans un intervalle unitaire autour de $P$. Alors $\kappa(P) = k\,\rho(P)$ est la baisse supplémentaire provoquée,
par les liquidations, pour un point de baisse au voisinage de $P$. Lorsque $\kappa$ est constant et inférieur à un, une baisse de $\delta$ se termine
en une baisse de $\delta/(1 - \kappa)$~; lorsque $\kappa \ge 1$, la cascade ne s'arrête pas tant que le prix ne quitte pas la zone.
Le prix final ne dépend ni de l'ordre ni de la taille des lots dans lesquels les positions déclenchées sont liquidées.
\end{proposition}

\begin{proof}
Une baisse $\delta$ déclenche $\rho\delta$ unités, lesquelles provoquent une baisse $\kappa\delta$, qui déclenche $\kappa^2\delta$, et ainsi
de suite~: la série géométrique a pour somme $\delta/(1 - \kappa)$ si $\kappa < 1$, et diverge sinon. Avec un impact linéaire et permanent,
le prix après un ensemble quelconque de liquidations est le prix choqué diminué de $k$ fois les unités vendues, quel que soit leur
ordre~; liquider chaque position dont le prix a été atteint, par lots quelconques, s'arrête au même plus grand point fixe.
\end{proof}

\begin{omfigure}
\begin{tikzpicture}
\begin{axis}[width=11cm,height=5cm,ybar,bar width=4pt,
  xlabel={prix (prix d'entrée autour de 100)},ylabel={$\kappa$},
  tick label style={font=\small},label style={font=\small},
  xmin=59,xmax=100,ymin=0,ymax=1.3,grid=major,grid style={gray!25},
  table/col sep=comma]
\addplot[fill=omThm!60,draw=omThm] table[x=price,y=kappa]{figdata/markets-3/18-margin-liquidation-and-loss-allocation/kappa.csv};
\draw[omDef,very thick,dashed] (axis cs:59,1) -- (axis cs:100,1);
\end{axis}
\end{tikzpicture}

\omcaption{Le multiplicateur local de cascade de la \cref{prop:m3:margin-liquidation-and-loss-allocation:kappa} pour les 4\,000 positions longues du tutoriel
(167~millions de dollars de notionnel, levier de 2 à 50 fois) et un impact de 0.12~\%
du prix par million de dollars vendu~; la ligne en pointillés marque $\kappa = 1$. Entre environ 92 et 96, le multiplicateur dépasse un~: une cascade qui atteint cette
zone la traverse. Positions synthétiques. Données~: le tutoriel du chapitre.}\label{fig:m3:margin-liquidation-and-loss-allocation:kappa}
\end{omfigure}

\Cref{fig:m3:margin-liquidation-and-loss-allocation:drawdown} montre la conséquence. De petits chocs sont amplifiés
de deux à trois fois, puis s'arrêtent. Un choc de 1.11~\% suffit à faire passer le prix dans la zone où $\kappa > 1$, et la
cascade descend alors sous 90~: une amplification par neuf. Plafonner le levier à dix fois supprime le groupe de \omterm{def:m3:margin-liquidation-and-loss-allocation:prices}{prix de liquidation}
proches du prix d'entrée et porte le seuil à 5.37~\%~; doubler la profondeur du carnet divise $\kappa$ par deux, et celui-ci reste inférieur à un
partout, de sorte que la baisse est amplifiée de façon progressive, à peu près par deux, et n'atteint 10~\% qu'après un choc de 5.55~\%~; sans
vente forcée, il n'y a aucune cascade. Le tutoriel vérifie que le résultat ne dépend pas de la taille des
lots du moteur, comme l'énonce la proposition, et désactive tour à tour chaque mécanisme.

\begin{omfigure}
\begin{tikzpicture}
\begin{axis}[width=11cm,height=5cm,
  xlabel={choc initial, \%},ylabel={baisse totale, \%},
  tick label style={font=\small},label style={font=\small},
  xmin=0,xmax=10,ymin=0,ymax=26,grid=major,grid style={gray!25},
  legend columns=2,legend style={font=\footnotesize,at={(0.5,-0.3)},anchor=north,draw=none,fill=none,/tikz/every even column/.append style={column sep=8pt}},
  table/col sep=comma]
\addplot[omThm,very thick] table[x=shock,y=base]{figdata/markets-3/18-margin-liquidation-and-loss-allocation/drawdown.csv};\addlegendentry{cas de base}
\addplot[omDef,very thick,dashed] table[x=shock,y=cap10]{figdata/markets-3/18-margin-liquidation-and-loss-allocation/drawdown.csv};\addlegendentry{levier plafonné à 10~fois}
\addplot[omExo,very thick,dotted] table[x=shock,y=depth2]{figdata/markets-3/18-margin-liquidation-and-loss-allocation/drawdown.csv};\addlegendentry{carnet deux fois plus profond}
\addplot[gray,thick] table[x=shock,y=noforced]{figdata/markets-3/18-margin-liquidation-and-loss-allocation/drawdown.csv};\addlegendentry{aucune vente forcée}
\end{axis}
\end{tikzpicture}

\omcaption{Chute totale après la cascade en fonction du choc initial, pour le cas de base et trois variantes. Dans
le cas de base, la baisse passe de moins de 4~\% à plus de 10~\% entre des chocs de 1.0~\% et de 1.25~\%~; plafonner le levier à
dix fois déplace le saut au-delà de 5~\%, et un carnet deux fois plus profond le fait disparaître. Positions et carnet synthétiques. Données~: le
tutoriel du chapitre.}\label{fig:m3:margin-liquidation-and-loss-allocation:drawdown}
\end{omfigure}

Un gap est pire qu'une glissade. Lorsque le cours décroche de 5~\% d'un coup, comme lors d'une annonce surprise, chaque position dont
le \omterm{def:m3:margin-liquidation-and-loss-allocation:prices}{prix de liquidation} se situe dans le gap est clôturée au prix du gap, sous de nombreux \omterm{def:m3:margin-liquidation-and-loss-allocation:prices}{prix de faillite}~: dans le tutoriel,
3.21~millions de dollars de déficits que le \omterm{def:m3:margin-liquidation-and-loss-allocation:fund}{fonds d'assurance} doit payer, et le cours finit près de 84. Le problème du week-end suit la
perte le long de la cascade de défaillance.

\section{Tutoriel~: prix de liquidation et cascade}\label{tut:m3:margin-liquidation-and-loss-allocation:cascade}

\begin{tutorial}
\textbf{Objectif.} Dériver les \omterm{def:m3:margin-liquidation-and-loss-allocation:prices}{prix de liquidation}, simuler une cascade de positions longues à effet de levier vendant dans un carnet de profondeur annoncée, tester
sa stabilité sous un pas de temps plus fin, et désactiver tour à tour chaque mécanisme. \textbf{État final~:}
les \cref{fig:m3:margin-liquidation-and-loss-allocation:liqprice,fig:m3:margin-liquidation-and-loss-allocation:kappa,fig:m3:margin-liquidation-and-loss-allocation:drawdown}
et les chiffres du texte.
\begin{tutsteps}
\item \textbf{\omterm{def:m3:margin-liquidation-and-loss-allocation:prices}{Prix de liquidation} et de faillite.} La \Cref{prop:m3:margin-liquidation-and-loss-allocation:liq} en code.
\omcode{code/firm/liquidation/firm_liquidation.py}{26}{37}{Prix de liquidation et de faillite d'une position linéaire isolée.}\label{lst:m3:margin-liquidation-and-loss-allocation:liq}
\item \textbf{La cascade.} Les positions déclenchées sont vendues dans le carnet, d'un seul coup ou par lots.
\omcode{code/markets-3/18-margin-liquidation-and-loss-allocation/python/m3_liq.py}{59}{83}{Une cascade de liquidations avec une taille de lot paramétrable.}\label{lst:m3:margin-liquidation-and-loss-allocation:cascade}
\item \textbf{Stabilité et désactivations.} Lancer \texttt{cascade} avec \texttt{batch} valant 100, 10 et 1, et comparer avec le
lot regroupant toutes les positions déclenchées~; puis poser \texttt{forced\_selling=False}, plafonner le levier et diviser l'impact par deux.
\item \textbf{Exécuter} \texttt{fig\_liq.py} pour les données du graphique.
\end{tutsteps}
\textbf{Pour aller plus loin.} Remplacer l'impact linéaire par un impact en racine carrée et vérifier si la taille des lots
reste sans effet~; laisser le carnet se reconstituer entre les lots~; ajouter des positions courtes et un paiement de funding.
\end{tutorial}

\section{Construction~: le moteur de marge et de liquidation}\label{bld:m3:margin-liquidation-and-loss-allocation:liquidation}

\begin{build}
\bfield{Rôle} La société miniature doit savoir, pour chaque position sur chaque plateforme, à quelle distance elle se trouve de la liquidation et
ce qu'une liquidation coûterait, et ses simulations de risque exigent que le moteur de la plateforme et sa cascade de défaillance soient reproduits.
\bfield{Interface} \texttt{Tier(max\_notional, mmr, maint\_amount)}, \texttt{maintenance\_margin(notional, tiers)}~;
\texttt{liq\_price\_linear}, \texttt{bankruptcy\_price\_linear}, \texttt{liq\_price\_inverse}~;
\texttt{cross\_health(wallet, positions, tiers)}~; \texttt{sell\_into\_book(qty, bids)}~; \texttt{adl\_rank(accounts,
price)}~; \texttt{liquidate(account, bids, fund, opposite, price)}.
\bfield{Règles} Les \omterm{def:m3:perpetual-futures:index}{prix de marquage} proviennent de \texttt{firm.perp}~; les paliers sont continus à leurs bornes~; l'excédent d'une liquidation revient au
fonds, et son déficit est payé par celui-ci~; si le fonds ne peut pas payer, la position est clôturée contre des positions opposées dans l'ordre
d'\omterm{def:m3:margin-liquidation-and-loss-allocation:adl}{auto-deleveraging} (ADL) au \omterm{def:m3:margin-liquidation-and-loss-allocation:prices}{prix de faillite}, le carnet et le fonds restant inchangés.
\bfield{Tests d'acceptation} \texttt{code/firm/liquidation/tests/}~: continuité des paliers~; \omterm{def:m3:margin-liquidation-and-loss-allocation:prices}{prix de liquidation} linéaires et
inverses confrontés aux formes closes~; santé de la \omterm{def:m3:margin-liquidation-and-loss-allocation:modes}{marge croisée}~; vente dans le carnet~; excédent versé au fonds~; ADL dans l'ordre de rang.
\bfield{Pour aller plus loin} Liquidation partielle vers un palier inférieur~; fournisseurs de liquidité de backstop~; une règle d'ADL au prorata~; exposition inter-plateformes
à un seul sous-jacent.
\end{build}

\begin{omsources}
\item CoinDesk Research, ``Market Spotlight: Inside Crypto's \$19 Billion Liquidation Event'', 17 octobre 2025.
\item Binance, ``What Is Auto-Deleveraging (ADL)'' et ``Leverage and Margin of USDS-M Futures'', pages d'aide~;
documentation Hyperliquid, ``Auto-deleveraging''. Consulté en septembre 2026.
\item BitMEX, points d'accès publics de l'API d'assurance et d'instruments, 23--24 septembre 2026~; blog BitMEX, ``How We Are Responding to the 13 March
DDoS Attacks'', mars 2020.
\item FTX, centre d'aide, ``Backstop Liquidity Provider Program'' (mis à jour le 14 septembre 2022), copie archivée sur Internet Archive.
\item Protos, ``Outdated algorithm caused \$650M excess losses on Hyperliquid, report'', 10 décembre 2025.
\end{omsources}

\section{Exercices}

\begin{exercise}[$\star$]\label{exo:m3:margin-liquidation-and-loss-allocation:1}
Calculer les \omterm{def:m3:margin-liquidation-and-loss-allocation:prices}{prix de liquidation} et de faillite d'une position longue à 20 fois, ouverte à 100, avec une marge de maintenance de 0.5~\%.
\end{exercise}

\begin{exercise}[$\star$]\label{exo:m3:margin-liquidation-and-loss-allocation:2}
Dire si un trader détenant une position à 50 fois le levier en \omterm{def:m3:margin-liquidation-and-loss-allocation:modes}{marge croisée}, avec un solde disponible important, est liquidé 2~\% sous
le prix d'entrée. Expliquer.
\end{exercise}

\begin{exercise}[$\star$]\label{exo:m3:margin-liquidation-and-loss-allocation:3}
Nommer les couches de la cascade de défaillance d'une plateforme de crypto-actifs, dans l'ordre.
\end{exercise}

\begin{exercise}[$\star\star$]\label{exo:m3:margin-liquidation-and-loss-allocation:4}
Comparer les \omterm{def:m3:margin-liquidation-and-loss-allocation:prices}{prix de liquidation} d'une position longue linéaire et d'une position longue inverse à 5 fois le levier et une marge de maintenance de 0.5~\%. Expliquer
pourquoi ils diffèrent.
\end{exercise}

\begin{exercise}[$\star\star$]\label{exo:m3:margin-liquidation-and-loss-allocation:5}
Près du prix courant $\kappa = 0.8$. Un choc d'actualité fait baisser le prix de 1~\%. Calculer la baisse totale si $\kappa$
reste constant.
\end{exercise}

\begin{exercise}[$\star\star$]\label{exo:m3:margin-liquidation-and-loss-allocation:6}
Un teneur de marché est long au comptant et vendeur du \omterm{def:m3:perpetual-futures:perp}{contrat perpétuel}. Expliquer comment l'\omterm{def:m3:margin-liquidation-and-loss-allocation:adl}{auto-deleveraging} peut lui nuire, et ce qu'il pourrait faire pour y remédier.
\end{exercise}

\begin{exercise}[$\star\star\star$]\label{exo:m3:margin-liquidation-and-loss-allocation:7}
\emph{Programmation.} Avec \texttt{critical\_shock}, déterminer le seuil pour des plafonds de levier de 20 et de 10 fois. Expliquer
ce que cela suggère sur les limites de levier imposées par les plateformes.
\end{exercise}

\begin{exercise}[$\star\star\star$]\label{exo:m3:margin-liquidation-and-loss-allocation:8}
\emph{Trouver l'erreur.} «~Une simulation de cascade qui liquide les positions une à une est plus précise qu'une simulation qui
liquide d'un seul coup toutes les positions déclenchées, donc sa réponse sera différente.~»
\end{exercise}

\section{Problème~: anatomie d'une cascade}

\begin{problem}[{Problème du week-end --- d'un choc à l'auto-deleveraging}]\label{pb:m3:margin-liquidation-and-loss-allocation:1}
Utiliser la population du tutoriel, soit 4\,000 positions longues à effet de levier (notionnel de 166.6~millions de dollars), un taux de marge de maintenance de 0.5~\%,
un impact de 0.0012 point de prix par unité vendue (0.12~\% pour 1~million de dollars à un prix de 100) et un \omterm{def:m3:margin-liquidation-and-loss-allocation:fund}{fonds d'assurance}
de 2~millions de dollars.

\medskip\noindent\textbf{Partie I --- Une position.}
\begin{enumerate}
\item Calculer la marge, le \omterm{def:m3:margin-liquidation-and-loss-allocation:prices}{prix de liquidation} et le \omterm{def:m3:margin-liquidation-and-loss-allocation:prices}{prix de faillite} d'une position longue de 250 unités à 100 avec un levier de 25 fois.
\item Si elle est liquidée à 95.5, que paie le \omterm{def:m3:margin-liquidation-and-loss-allocation:fund}{fonds d'assurance}~?
\item Et à 96.8~?
\item Pourquoi le fonds reçoit-il de l'argent dans le second cas~?
\item Pourquoi la plateforme utilise-t-elle le \omterm{def:m3:perpetual-futures:index}{prix de marquage}, et non la dernière transaction, pour la déclencher~?
\end{enumerate}

\medskip\noindent\textbf{Partie II --- La cascade.}
\begin{enumerate}[resume]
\item Quel est le plus petit choc initial après lequel le prix baisse de 10~\% ou plus au total~?
\item Que dépasse $\kappa$ entre environ 92 et 96, et que cela signifie-t-il~?
\item Quels sont les seuils lorsque le levier est plafonné à 10 fois, et lorsque le carnet est deux fois plus profond~?
\item Pourquoi le prix final est-il le même quelle que soit la taille du lot~?
\item Qu'est-ce qui remettrait en cause cette invariance~?
\end{enumerate}

\medskip\noindent\textbf{Partie III --- La cascade de défaillance.}
\begin{enumerate}[resume]
\item Après un gap de 5~\%, où finit le prix, et quels sont les déficits~?
\item Quelle part atteint l'\omterm{def:m3:margin-liquidation-and-loss-allocation:adl}{auto-deleveraging}~?
\item Qui est désendetté en premier selon un classement par pourcentage de P\&L multiplié par le levier~?
\item Que perd un teneur de marché désendetté, et que conserve-t-il~?
\item Comment une règle au prorata modifierait-elle la répartition de la charge~?
\end{enumerate}

\medskip\noindent\textbf{Partie IV --- Jugement.}
\begin{enumerate}[resume]
\item Les plateformes doivent-elles plafonner le levier~?
\item Un \omterm{def:m3:margin-liquidation-and-loss-allocation:fund}{fonds d'assurance} de 2~millions de dollars suffit-il pour ce marché~?
\item Comment une société de trading doit-elle utiliser l'indicateur ADL d'une plateforme~?
\item Énoncer le \emph{résultat nommé}~: le choc critique, et la perte qui atteint le \omterm{def:m3:margin-liquidation-and-loss-allocation:fund}{fonds d'assurance} puis
l'\omterm{def:m3:margin-liquidation-and-loss-allocation:adl}{auto-deleveraging} après un gap de 5~\%.
\item En une phrase~: qui se trouve derrière une transaction de produits dérivés sur crypto-actifs~?
\end{enumerate}
\end{problem}

\section{Questions d'entretien}

\begin{interviewq}[$\star$ \iqroles{trader}]\label{iq:m3:margin-liquidation-and-loss-allocation:1}
Calculez le \omterm{def:m3:margin-liquidation-and-loss-allocation:prices}{prix de liquidation} d'une position longue à 10 fois.
\end{interviewq}

\begin{interviewq}[$\star$ \iqroles{risque}]\label{iq:m3:margin-liquidation-and-loss-allocation:2}
Qu'est-ce que l'\omterm{def:m3:margin-liquidation-and-loss-allocation:adl}{auto-deleveraging}, et pourquoi une société qui n'est pas endettée devrait-elle s'en préoccuper~?
\end{interviewq}

\begin{interviewq}[$\star\star$ \iqroles{chercheur}]\label{iq:m3:margin-liquidation-and-loss-allocation:3}
Qu'est-ce qui détermine si un choc de prix se transforme en \omterm{def:m3:margin-liquidation-and-loss-allocation:cascade}{cascade de liquidations}~?
\end{interviewq}

\begin{interviewq}[$\star\star$ \iqroles{développeur}]\label{iq:m3:margin-liquidation-and-loss-allocation:4}
Comment concevriez-vous un \omterm{def:m3:margin-liquidation-and-loss-allocation:engine}{moteur de liquidation} pour minimiser les pertes subies par le \omterm{def:m3:margin-liquidation-and-loss-allocation:fund}{fonds d'assurance}~?
\end{interviewq}

\begin{interviewq}[$\star\star$ \iqroles{risque, trader}]\label{iq:m3:margin-liquidation-and-loss-allocation:5}
Comparez \omterm{def:m3:margin-liquidation-and-loss-allocation:modes}{marge isolée} et \omterm{def:m3:margin-liquidation-and-loss-allocation:modes}{marge croisée} pour un teneur de marché cotant 50 \omterm{def:m3:perpetual-futures:perp}{contrats perpétuels}.
\end{interviewq}

\begin{interviewq}[$\star\star\star$ \iqroles{chercheur, risque}]\label{iq:m3:margin-liquidation-and-loss-allocation:6}
À partir de données publiques, comment estimeriez-vous où se situent les \omterm{def:m3:margin-liquidation-and-loss-allocation:prices}{prix de liquidation} de l'encours ouvert d'une plateforme~?
\end{interviewq}
